\documentclass[9pt,a4paper,sans]{moderncv}
\moderncvstyle{casual}
\moderncvcolor{blue}
\usepackage[scale=0.75]{geometry}
\usepackage[T1]{fontenc}
\usepackage[utf8x]{inputenc}
\usepackage[spanish]{babel}
\usepackage{bibunits}
\usepackage{xcolor}

\newcommand{\Q}{{\textsf{Q}\hspace*{-1.1ex}%
  \rule{0.15ex}{1.5ex}\hspace*{1.1ex}}}
\newcommand{\Cuat}{º\Q~}
\newcommand{\actual}{$\infty$}
\newcommand{\vat}{CUIT}
\newcommand{\birthdate}{Fecha de Nacimiento}

\firstname{Cristian Sebastian}
\familyname{Rocha}
\title{Curriculum Vitae}
\address{Lavalle 2294 2°}{Villa Ballester, Gral. San Martin, Buenos Aires (1653)}
\phone[mobile]{(+54-911)~6800~0269}
\email{csrocha@gmail.com}
\homepage{https://github.com/csrocha}
\extrainfo{\vat 23-25095454-9 - \birthdate: 4 dic 1975}

\begin{document}
\makecvtitle

\section{Resumen Profesional}
Líder tecnológico con más de 20 años de experiencia en el sector IT, especializado en la **automatización de procesos complejos e implementación de soluciones de Inteligencia Artificial**. Experto en el diseño de infraestructura para procesamiento de datos a gran escala y en la industrialización de modelos analíticos. Cuento con una sólida trayectoria liderando equipos multidisciplinarios, definiendo estrategias tecnológicas y transformando requerimientos de negocio en productos innovadores de alto impacto.

\section{Habilidades Principales}
\cvitem{Estrategia e IA}{Liderazgo de proyectos de IA Generativa, agentes inteligentes para curación de datos, automatización de punta a punta (End-to-End) y definición de roadmaps tecnológicos.}
\cvitem{Automatización}{Software Reliability Engineering (SRE), industrialización de procesos estadísticos, CI/CD, y desarrollo de herramientas de automatización de infraestructura.}
\cvitem{Stack Técnico}{Python, R, C++, JavaScript, SQL/NoSQL (Cassandra), Kafka, infraestructura en la nube, y plataformas de IA/ML.}
\cvitem{Liderazgo}{Gestión de equipos de ingeniería, mentoría técnica, negociación con proveedores y alineación de tecnología con objetivos de negocio.}

\section{Experiencia Profesional}

\cventry{jul 2022--Presente}
	{Jefe de Sistema (CTO)}
	{Fundación Observatorio PyME}
	{Buenos Aires, Argentina}
	{}
	{Liderazgo de la transformación digital y la estrategia de IA de la organización.
	\begin{description}
		\item[Industrialización de Datos:] Automatización integral de procesos internos para la generación de estadísticas y reportes. Desarrollo de agentes de IA para detección de fallas en microdatos.
		\item[IA Generativa:] Implementación de una plataforma de scoring para PyMEs utilizando IA generativa para la creación automatizada de más de 3000 informes explicativos de resultados.
		\item[Estrategia y Liderazgo:] Dirección de equipos de desarrollo e infraestructura, asegurando la escalabilidad de los productos y la seguridad de los datos sensibles.
	\end{description}}

\cventry{abr 2020--jul 2022}
	{Ingeniero de Infraestructura para IA y Procesamiento de Datos}
	{Mercado Libre}
	{Buenos Aires, Argentina}
	{}
	{Responsable del diseño y optimización de la plataforma centralizada de IA.
	\begin{description}
		\item[Plataforma de IA:] Diseño y desarrollo de la infraestructura común utilizada por todos los equipos de científicos de datos y desarrolladores de IA de la compañía.
		\item[Optimización:] Mejora del rendimiento y eficiencia en el uso de recursos de cómputo para modelos de Machine Learning.
		\item[Seguridad y Resiliencia:] Adaptación de la plataforma a normativas internacionales de seguridad de datos y mejora de la tolerancia a fallos.
	\end{description}}

\cventry{jul 2016--mar 2020}
	{Associated - Software Reliability Engineer (SRE)}
	{JPMorgan Chase}
	{Buenos Aires, Argentina}
	{}
	{Enfoque en la automatización de infraestructura y confiabilidad de sistemas críticos.
	\begin{description}
		\item[Automatización:] Diseño y programación de aplicaciones para la automatización de infraestructura utilizando Python, Bash, Kafka y Cassandra.
		\item[SRE:] Implementación de prácticas de ingeniería de confiabilidad para asegurar la disponibilidad de aplicaciones bancarias globales.
	\end{description}}

\cventry{2009--2016}
	{Consultor Senior y Emprendedor}
	{Cooperativa Moldeo Interactive / Proyectos Varios}
	{Buenos Aires, Argentina}
	{}
	{Liderazgo técnico en proyectos de automatización y análisis de datos.
	\begin{description}
		\item[SAS ABT Builder:] Desarrollador principal de un sistema de generación automática de tablas ABT para grandes entidades financieras (Telecom, Banco Itau, Banco Supervielle).
		\item[GeoStat:] Diseño de una plataforma de información geoestadística para PyMEs utilizando R, Python y D3JS.
		\item[Consultoría SAS:] Especialista en Marketing Automation y soluciones de Business Intelligence para el sector bancario.
	\end{description}}

\cventry{2014--2016}
	{Arquitecto de Datos y Desarrollador Full Stack}
	{Ministerio Público Fiscal - Fundación Sadosky}
	{Buenos Aires, Argentina}
	{}
	{Diseño e implementación de la Base de Datos de Conocimientos para la Unidad Especializada en casos de apropiación de niños. Desarrollo de módulos de seguimiento, relaciones entre personas y generación de estadísticas sobre plataforma Odoo.}

\section{Formación Académica}
\cventry{2002}{Licenciado en Ciencias de la Computación}{Universidad de Buenos Aires (UBA)}{Argentina}{}{}
\cventry{2004--2012}{Doctorado en Ciencias de la Computación (Candidato)}{UBA}{Argentina}{}{Investigación enfocada en algoritmos para predicción de estructuras proteicas y bioinformática.}
\cventry{2001}{Trainer in Bioinformatics}{Oswaldo Cruz Foundation / WHO}{Brasil}{}{Especialización en computación aplicada a genómica.}

\section{Trayectoria Académica y Científica}
\cvitem{Docencia}{Jefe de Trabajos Prácticos en la UBA (Algoritmos II, Organización del Computador). Más de 10 años formando profesionales en ingeniería de software y computación científica.}
\cvitem{Investigación}{Investigador en Biología Computacional (UBA, ANSES, Hospital Italiano). Desarrollo de algoritmos de alto rendimiento en C++ y CUDA para el procesamiento de información biológica.}
\cvitem{Publicaciones}{Autor de diversas publicaciones científicas y director de tesis de licenciatura enfocadas en algoritmos y optimización.}

\section{Idiomas}
\cvlistdoubleitem{Español: Nativo}{Inglés: Avanzado}

\end{document}
